% 图打散 = 以参考格 f 为中心，把网格分成四个区域，每个区域都要有可实现模式。
% 区域（相对 f=(f1,f2)）：f 所在格 = 两处都匹配；f 所在行其余格 = 只匹配 x1；
% f 所在列其余格 = 只匹配 x2；其余四格 = 两处都不匹配。
% 左：f=(正常,广告)。见证：h2(正常,广告)、h1(正常,正常)、h3(广告,广告)、h5(诈骗,诈骗)。
%     四区都有实心点，{x1,x2} 被图打散，d_G(V)=2。“两处都不匹配”的见证取 h5，
%     它在 x1 上偏离为诈骗、在 x2 上偏离为诈骗；h3、h1 的偏离标签分别是广告、正常，替代标签不固定。
% 右：f=(正常,诈骗)。该格不在 V 中，“两处都匹配”无法实现，此参考失败；
%     图打散只要求存在某个参考 f。
\begingroup
\definecolor{LtSixGreen}{HTML}{829C8B}
\definecolor{LtSixGreenFill}{HTML}{E8EEE9}
\begin{tikzpicture}[figure picture,x=1mm,y=1mm]
 \def\cs{10.5}
 \newcommand{\ltsixcell}[4]{% #1 off, #2 row, #3 col, #4 fill
   \fill[#4] ({#1+\cs*(#3-1)},{21-\cs*(#2-1)}) rectangle ++(\cs,\cs);}
 \newcommand{\ltsixframe}[2]{%
  \foreach \r in {1,2,3}{\foreach \c in {1,2,3}{
    \draw[FigureGrid,line width=.5pt] ({#1+\cs*(\c-1)},{21-\cs*(\r-1)}) rectangle ++(\cs,\cs);}}
  \foreach \lab [count=\c] in {正常,广告,诈骗}{
    \node[anchor=south,font=\figurefont\scriptsize] at ({#1+\cs*(\c-.5)},32) {\lab};}
  \foreach \lab [count=\r] in {正常,广告,诈骗}{
    \node[anchor=east,font=\figurefont\scriptsize] at ({#1-1},{26.25-\cs*(\r-1)}) {\lab};}
  \node[figure heading] at ({#1+1.5*\cs},40) {#2};
  \foreach \r/\c in {1/1,1/2,2/2,2/3,3/3,3/1}{
    \fill[FigureBlue] ({#1+\cs*(\c-.5)},{26.25-\cs*(\r-1)}) circle (1.6);}
  \foreach \r/\c in {1/3,2/1,3/2}{
    \draw[FigureMuted,line width=.6pt,fill=white] ({#1+\cs*(\c-.5)},{26.25-\cs*(\r-1)}) circle (1.2);}
 }
 % ---- left: f=(正常,广告) at row1 col2
 \ltsixcell{9}{1}{2}{FigureYellowFill}
 \ltsixcell{9}{1}{1}{FigureBlueFill}\ltsixcell{9}{1}{3}{FigureBlueFill}
 \ltsixcell{9}{2}{2}{LtSixGreenFill}\ltsixcell{9}{3}{2}{LtSixGreenFill}
 \foreach \r/\c in {2/1,2/3,3/1,3/3}{\ltsixcell{9}{\r}{\c}{FigureRedFill}}
 \ltsixframe{9}{参考 $f=(\text{正常},\text{广告})$：成功}
 \foreach \r/\c in {1/2,1/1,2/2,3/3}{
   \draw[FigureInk,line width=.7pt] ({9+\cs*(\c-.5)},{26.25-\cs*(\r-1)}) circle (2.7);}
 % ---- right: f=(正常,诈骗) at row1 col3
 \ltsixcell{60}{1}{3}{FigureYellowFill}
 \ltsixcell{60}{1}{1}{FigureBlueFill}\ltsixcell{60}{1}{2}{FigureBlueFill}
 \ltsixcell{60}{2}{3}{LtSixGreenFill}\ltsixcell{60}{3}{3}{LtSixGreenFill}
 \foreach \r/\c in {2/1,2/2,3/1,3/2}{\ltsixcell{60}{\r}{\c}{FigureRedFill}}
 \ltsixframe{60}{参考 $f=(\text{正常},\text{诈骗})$：失败}
 \draw[FigureRed,line width=.85pt] ({60+\cs*2.5-2.2},{26.25-2.2})--++(4.4,4.4)
   ({60+\cs*2.5-2.2},{26.25+2.2})--++(4.4,-4.4);
 % ---- legend
 \foreach \fillc/\txt [count=\i] in {FigureYellowFill/两处都匹配,FigureBlueFill/只匹配$x_1$,
     LtSixGreenFill/只匹配$x_2$,FigureRedFill/两处都不匹配}{
   \fill[\fillc] ({4+23*(\i-1)},-7) rectangle ++(4,3.2);
   \draw[FigureGrid,line width=.4pt] ({4+23*(\i-1)},-7) rectangle ++(4,3.2);
   \node[anchor=west,figure note] at ({9+23*(\i-1)},-5.4) {\txt};}
 \draw[FigureInk,line width=.7pt] (99,-5.4) circle (1.6);
 \node[anchor=west,figure note] at (101.5,-5.4) {见证};
\end{tikzpicture}
\endgroup
